\section{创新点一}
创新点1：动态视觉启发的极限场景感知。
提出了面向高速目标的、面向微小目标和面向暗光目标的生物启发式算法，解决了城市交通场景中高速运动目标难以稳定跟踪、城市工厂场景中微小缺陷与细粒度作业目标难以精准识别、城市安防与城市应急场景中弱小目标和暗光目标难以可靠感知等难题，显著提升了复杂城市场景下高速微小目标的全时感知能力，检测时延最低可至3毫秒，较同类方法提升20%以上；达到亚像素级检测精度，可在55米外分辨0.33毫米细线；暗光（<5 lux）条件下重建SSIM均值>0.75，提升20%以上。

1.1 事件触发的高速时空响应与毫秒级目标定位
针对城市交通与低空飞行场景中高速运动目标易产生运动模糊、帧间状态信息缺失，导致目标难以连续跟踪、稳定定位和快速响应的问题，提出了事件触发的高速时空响应与目标检测定位方法。通过异步事件流高频运动信息提取、事件—深度协同跟踪和事件—雷达时空互补融合，连续恢复传统帧相机采样间隔内的目标位置、姿态与运动状态，实现了高速目标的毫秒级连续跟踪、三维定位，和实时避障。代表性系统目标检测时延最低可至3毫秒，较同类方法提升20%以上，并使无人机跟踪精度提升22%以上；无人机平均定位端到端时延约5ms，较同类方法提升60%以上，为高速动态目标毫秒级感知与定位提供了量化支撑。
直接支撑论文：EventPro，mmE-Loc

1.2 事件边缘空间驱动的精细表征与亚像素级识别
针对城市工业场景中微小缺陷和细粒度作业目标像素占比低、边缘特征弱，以及图像、事件和点云等异构信息难以精确对齐，导致微小目标细节难表征、空间位置难解算的问题，提出了事件边缘空间驱动的跨模态精细表征与高精度解算方法。通过将图像纹理、事件变化和点云几何信息统一编码至事件边缘空间，构建跨模态显式对齐、可靠性感知融合和二维—三维运动联合约束机制，增强微小目标的轮廓结构与细粒度运动特征，实现了复杂退化条件下的精细目标识别和高精度空间定位。代表性系统可达到亚像素级检测精度，可在55米外分辨0.33毫米细线，且检测准确率96.7%。
直接支撑论文：x²-Fusion，SkyShield

1.3 双时间尺度脉冲协同的极端环境增强与暗光感知
针对城市安防与应急场景中弱小目标有效信号稀疏、暗光目标对比度低，以及强光和复杂背景干扰导致目标信息易被噪声淹没的问题，提出了双时间尺度脉冲协同的极端环境视觉增强方法。通过短时间尺度提取目标瞬态变化、长时间尺度聚合稳定结构信息，并结合四维事件云建模、前景—背景事件分离和弱光多模态信息增强，抑制环境噪声并累积微弱目标响应，实现了强干扰和低照度条件下事件流在线净化、弱小目标增强与暗光目标可靠感知。代表性系统在低于5 lux条件下，环境重建结果的SSIM均值大于0.75，较同类方法提升>20%；复杂天气感知去噪的信号保持率达 0.95、噪声去除率达 0.91；模型仅 0.26M 参数，推理速度相比同类型方法提升约 25 倍。
直接支撑论文：PRE-Mamba
